% Source: 03 - Wed Sep 30 2026.pdf, page 13. Content remains in source order.
\sourcepage{03}{13}
Fortunately, $\polyred$ is transitive!

If I know that a given language $\hat L\in\classNPC$, it suffices to prove that
\[
\underbrace{\hat L}_{\text{Known to be in NPC}}\polyred
\underbrace{\bar L}_{\text{Candidate to be in NPC}}.
\]
If $\hat L\polyred\bar L$, since $\forall L\in\classNP:\ L\polyred\hat L$ (because $\hat L\in\classNPC$) we have that for every language in NP, we can reduce it to $\hat L$ in polinomial time, and then reduce $\hat L$ to $\bar L$ in polynomial time. By transitivity of $\polyred$, we have that for every language in NP, we can reduce it to $\bar L$ in polynomial time. So $\bar L$ is NP-hard, and if we can also show that $\bar L\in\classNP$, then we have shown that $\bar L\in\classNPC$.
\[
\forall L\in\classNP:\ (L\polyred\hat L)\land(\hat L\polyred\bar L)
\xRightarrow{\text{transitivity of }\polyred}\forall L\in\classNP:\ L\polyred\bar L.
\]
A single reduction suffices!

Unfortunately, the transitivity trick cannot be used for the first language/problem to be placed in NPC!
The first problem requires a ``universal'' reduction from a generic problem in NP, proving all languages in NP can be reduced to it. Computer scientist Cook was the first to show that the SAT problem is NP-complete, by showing that every language in NP can be reduced to SAT in polynomial time. This is known as Cook's Theorem.